\subsection{Evaluatie}
\label{sec:evaluatie}

Om de bruikbaarheid van het algoritme te evalueren, wordt gebruikgemaakt van de testset die tijdens de voorbereiding van de data werd afgesplitst (sectie~\ref{subsec:voorbereiding-data}). De resultaten van het model worden gerapporteerd aan de hand van de metrieken uit sectie~\ref{subsubsec:evaluatiemetrieken}, met het false-positive-percentage (\emph{FPR}) als de belangrijkste. Deze metriek geeft immers aan hoe vaak het algoritme een legitieme e-mail ten onrechte als phishing aanduidt. Zoals besproken in sectie~\ref{subsubsec:uitleg-voor-gebruikers}, besteden gebruikers steeds minder aandacht aan waarschuwingen die ze herhaaldelijk te zien krijgen \autocite{Anderson2015}. Een hoge \emph{FPR} brengt dan ook het risico mee dat gebruikers aan de waarschuwingen van het algoritme wennen, waardoor dit onderzoek zijn doel mist.
